% Generated by Claude Fable 5.1 (Anthropic), 2026-09-15 — Section III "Extraction" (v3, layer numbers as in Fig. 2: layer 3 = entity, layer 4 = callback)
\noindent\textbf{Extraction.} Initialization events yield the vertices of
Fig.~\ref{fig:model}a and the message and service edges the source code
declares via API. \texttt{sub} and \texttt{srv} edges are first created at
the entity layer, between interfaces: the \texttt{rcl\_*\_init} events give
each interface, its handle and the topic or service it is bound to, and an
edge joins every publisher to every subscription on its topic, every client
to its server. The \texttt{rclcpp\_*\_callback\_added} events supply the
containment below: the callback each subscription, service or timer
dispatches to. The communication events of the execution phase then let
these edges find their real owners at the callback layer: the publishing and
requesting side has no registration, so which callback issued a message or a
request is read from the FISH \texttt{publish\_link} and \texttt{client\_link}
events, fired inside the \texttt{callback\_start}--\texttt{callback\_end}
window of the issuer. Two further kinds have no declaration at all and are
obtained by observing patterns in the execution trace
(Fig.~\ref{fig:model}b). A \texttt{polled-msg} edge: node $Y$ subscribes to
topic $T$ through subscription $S$, but instead of $S$'s own callback,
another callback $R$ of $Y$ takes the message via \texttt{rmw\_take} during
its own execution. A \texttt{join}: in one node, a topic $T$ is published
from two or more different subscription callbacks, each time by the one that
ran last, the input that completed the set. With every edge in place at the
callback layer, each carrying one interaction, the graph is lifted upwards
(Fig.~\ref{fig:model}c): at the entity layer the lifted edges coincide with
the declared ones, except for declared pairs the run never exercised, which
exist only there; between nodes, and between executors, the edges among the
callbacks they contain are collected into one edge carrying the set of
interactions, while pairs that fall inside one vertex stay contained in it.
Each GPU-submitting callback receives the DAG, shapes and conditional graph
of Fig.~\ref{fig:model}d. Phases are delimited from the trace itself:
initialization ends when every periodic timer has fired once, warm-up when
every recurring callback has completed its first execution, and what follows
is the steady state.
